% Editorial ALC completion sheet for Part IX.
\calculussheet{sheet:tableaux-alc}{Blocking Completion Tableau for \(\mathcal{ALC}\)}{Data: an optional finite tested \(\mathcal{ALC}\) concept \(C_0\), a finite TBox \(\mathcal T\), and a finite ABox \(\mathcal A\), normalized as in Definition~\ref{def:p09-alc-tableau}; shape: a branching completion graph whose nodes \(x\) carry labels \(\mathcal L(x)\) and whose edges carry role names; declared ABox edges may join roots, while freshly generated existential edges form a successor forest; acceptance: one complete open branch, meaning that no rule remains applicable and no clash occurs; equality: none; blocking: a nonroot node \(y\) is directly blocked by an ancestor \(x\) in the generated successor forest when \(\mathcal L(y)\subseteq\mathcal L(x)\), and is indirectly blocked when an ancestor in that forest is blocked; every comparison is re-evaluated as labels grow.}
\noindent\emph{Editorial standing:} this exact sheet is the handbook's fixed
calculus with the proof identity in Annex~\ref{annex:p9.alc.1}; it is not a
transcription of either body-unverified registered book presentation.

As in Definition~\ref{def:p09-alc-tableau},
\(\mathsf{Glob}(\mathcal T)\) abbreviates the finite set of NNF concepts
\(\operatorname{nnf}(\neg C\sqcup D)\) for
\(C\sqsubseteq D\in\mathcal T\).  The abbreviation is metalinguistic, not a
new concept constructor; it states once the global condition used by the
initialization and successor rules below.

\hbrule{rule:p09-alc-initial}{Initialization}{\mathcal T,\mathcal A\quad\text{and optional }C_0}{\begin{gathered}\mathcal L(a)=\{\operatorname{nnf}(D):D(a)\in\mathcal A\}\cup\mathsf{Glob}(\mathcal T),\\r(a,b)\in\mathcal A\Rightarrow a\mathrel r b,\\\mathcal L(c_0)=\{\operatorname{nnf}(C_0)\}\cup\mathsf{Glob}(\mathcal T)\end{gathered}}{create one root for every ABox individual and retain every declared role edge; when \(C_0\) is present create the fresh root \(c_0\); if no root is otherwise present, create an anonymous root labelled \(\{\top\}\cup\mathsf{Glob}(\mathcal T)\)}
\hbrule{rule:p09-alc-and}{\(\sqcap\)}{x:C\sqcap D}{x:C,\ D}{add both conjuncts to \(\mathcal L(x)\) if either is absent}
\hbrule{rule:p09-alc-or}{\(\sqcup\)}{x:C\sqcup D}{x:C\ \mid\ x:D}{branch if neither disjunct is present; copy the current branch and add one displayed alternative to each child}
\hbrule{rule:p09-alc-exists}{\(\exists\)}{x:\exists r.C}{x\mathrel r y,\quad y:C,\quad y:\mathsf{Glob}(\mathcal T)}{\(x\) is unblocked and has no existing \(r\)-successor whose label contains \(C\); \(y\) is fresh on the branch}
\hbrule{rule:p09-alc-forall}{\(\forall\)}{x:\forall r.C\qquad x\mathrel r y}{y:C}{add \(C\) whenever it is absent; the rule remains applicable to every later \(r\)-successor, including successors created after the universal formula was first processed}
\hbrule{rule:p09-alc-clash}{clash}{x:\bot\quad\text{or}\quad x:A,\neg A}{\text{closed branch}}{\(A\) is a concept name; a closed branch receives no further expansion}

Only the existential rule is disabled at a blocked source.  Boolean and
universal consequences are still added at blocked nodes and across every edge
already present.  If label growth destroys a subset comparison, the node
becomes unblocked and its still-unsatisfied existential requirements become
applicable.  Rule applications that neither add a formula nor create a witness
are not legal expansion steps.
